%% ================================================================================================
%%  Appendix A. Proofs
%%  Code checks: simulations/python/voting.py --check (Propositions 4-7, the PVC test),
%%  experiments.py --check (Proposition 2), examples.py --check (Examples E1-E4)
%% ================================================================================================
\section{Proofs}\label{app:proofs}

\paragraph{Proposition~\ref{prop:majority-veto}.}
Take $n$ large enough that some integer $k$ satisfies $\theta n < k < n/2$, and options
$x, y, z_1, \dots, z_{m-2}$. Let $n-k$ voters rank $x \pref y \pref z_1 \pref \dots \pref z_{m-2}$ and
veto nothing, and $k$ voters rank $y \pref z_1 \pref \dots \pref z_{m-2} \pref x$ and veto $x$ only. Since
$n - k > n/2$, $x$ beats every other option: it is the Condorcet winner. It is vetoed by
$k > \theta n$ voters, while $y$ is vetoed by none. A rule that elects the Condorcet winner elects here an
option vetoed by at least $\theta n$ voters although $y$ is vetoed by fewer, which proves the first
claim. For the second, the coalition $T$ of the $k$ minority voters unanimously prefers the $m-1$ other
options to $x$; it blocks $x$ if $m - 1 \ge m - \lceil mk/n \rceil + 1$, that is, if
$\lceil mk/n \rceil \ge 2$, which holds when $k/n > 1/m$. \hfill$\square$

\paragraph{Proposition~\ref{prop:gridlock}.}
(i) Suppose an option $y \neq \sq$ is strictly preferred to $\sq$ by every voter:
$\|y - p_i\| < \|\sq - p_i\|$ for all $i$. Expanding the squares, this is
$(y - \sq) \cdot p_i > (\|y\|^2 - \|\sq\|^2)/2$ for all $i$: every ideal point lies in the open half-space
$H = \{z : (y - \sq) \cdot z > (\|y\|^2 - \|\sq\|^2)/2\}$. $H$ is convex, so it contains the convex hull of
the ideal points. But $\sq \notin H$, since $(y - \sq) \cdot \sq - (\|y\|^2 - \|\sq\|^2)/2 = -\|y-\sq\|^2/2 < 0$.
Hence $\sq$ is not in the convex hull, a contradiction. Under unanimity with the status quo as default,
a proposal passes only if every voter prefers it to $\sq$, so none passes.

(ii) Let $U$ be an open ball around $\sq$ contained in the support. Choose a simplex
$\Delta \subset U$ with vertices $a_0, \dots, a_d$ that contains $\sq$ in its interior, and $\varepsilon > 0$
small enough that every simplex with one vertex in each ball $B(a_j, \varepsilon)$ still contains $\sq$
(the barycentric coordinates of $\sq$ depend continuously on the vertices). Each ball $B(a_j,\varepsilon)$
lies in the support and has probability $q_j > 0$. If every ball contains at least one ideal point,
$\sq$ is in their convex hull. The probability that some ball is empty is at most
$\sum_j (1 - q_j)^n$, which tends to zero. \hfill$\square$

\paragraph{Proposition~\ref{prop:pvc-clones}.}
In Example~E1 ($n = 100$, $m = 3$), the 49 minority voters prefer $B$ and $C$ to $A$; since
$2 \ge 3 - \lceil 3 \cdot 49/100 \rceil + 1 = 2$, they block $A$, which is outside the core. In Example~E2
($m = 5$), they still prefer only $B$ and $C$ to $A$, and $2 < 5 - \lceil 5 \cdot 49/100 \rceil + 1 = 3$:
they no longer block $A$. No other coalition does either, since a coalition that contains a majority
voter has no option that all its members prefer to $A$, and a smaller coalition of minority voters has
less veto power. So $A$ enters the core, and as the Condorcet winner it is elected by Schulze within the
core, while $B$ was elected before (Table~\ref{tab:examples}). For fallback bargaining, in E1 the
majority's top two $\{A, B\}$ and the minority's $\{B, C\}$ meet at $B$ at depth 2. In E2 the majority's
top two are $\{A, A_2\}$ and nothing is common at depth 2; at depth 3 the majority's $\{A, A_2, A_3\}$ and
the minority's $\{B, C, A\}$ meet at $A$, which wins. \hfill$\square$

\paragraph{Proposition~\ref{prop:decisive}.}
The admissible set is never empty: either some option is vetoed by fewer than $\theta n$ voters, or the
fallback keeps the least-vetoed options. The Schulze method always has at least one winner on a non-empty
set \citep{Schulze2011}, and the tie-break ends with a strict order, so exactly one option is selected. All
steps use only the ballots, and treat every option, the status quo included, alike; the order of
proposal enters nowhere. \hfill$\square$

\paragraph{Proposition~\ref{prop:acceptability}.}
If some option has $\veto(x) < \theta n$, the admissible set is $\{x : \veto(x) < \theta n\}$ and the
winner belongs to it. Otherwise the winner belongs to $\arg\min \veto$. In both cases the winner is vetoed
by at least $\theta n$ voters only if every option is, which no rule can avoid. \hfill$\square$

\paragraph{Proposition~\ref{prop:condorcet-admissible}.}
The Schulze method is Smith-consistent, hence Condorcet-consistent \citep{Schulze2011}; applied to the
profile restricted to $\Adm$, its winners lie in the Smith set of $\Adm$, and the tie-break selects among
them. If the Condorcet winner $c$ of $\Opts$ has $\veto(c) < \theta n$, then $\Adm$ is the set of options
vetoed by fewer than $\theta n$ voters, $c \in \Adm$, and $c$ beats every option of $\Adm$, so it wins.
For Pareto, let $y$ dominate $x$: every voter weakly prefers $y$ to $x$, some strictly. By consistency of the
veto sets, any voter who vetoes $y$ also vetoes $x$, so $\veto(y) \le \veto(x)$ and $y \in \Adm$ whenever
$x \in \Adm$, in both branches of the definition. Within $\Adm$, $y$ beats $x$, and the Schulze method never
elects a Pareto-dominated option \citep{Schulze2011}. \hfill$\square$

\paragraph{Proposition~\ref{prop:qvc-clones}.}
Let $x'$ be a clone of $x$. By Definition~\ref{def:clones}, $\veto(x') = \veto(x)$ and the veto counts of
the other options are unchanged. In the first branch, $x' \in \Adm$ if and only if $x \in \Adm$, and the
other options keep their status; in the fallback branch, the minimum veto count is unchanged and
$x' \in \arg\min \veto$ if and only if $x$ is. So the admissible set either gains $x'$ (if it contains
$x$) or is unchanged, and $\{x, x'\}$ is a clone set in the profile restricted to it. The Schulze method is
independent of clones \citep{Schulze2011}; the first tie-break (fewest vetoes) gives clones equal counts,
and the second applies Schulze again. The winner is therefore unchanged up to the clone set, except when
the final random order separates exactly tied options. \hfill$\square$

\paragraph{Proposition~\ref{prop:blocking}.}
A coalition $C$ with $|C| \ge \theta n$ that vetoes $x$ raises $\veto(x)$ to at least $\theta n$, so that
$x \notin \Adm$ unless no option is vetoed by fewer than $\theta n$ voters. For the second claim, let $n$
be a multiple of $m$ and split the voters into $m$ equal blocs; bloc $j$ ranks the options in the cyclic
order $o_j \pref o_{j+1} \pref \dots \pref o_{j-1}$ (indices modulo $m$) and vetoes nothing, so every
option is admissible. Whatever the winner $w = o_k$, every bloc except bloc $k$ prefers $o_{k-1}$ to $w$:
a coalition of $(m-1)n/m \ge \theta n$ voters prefers another option and can make $w$ inadmissible by
vetoing it, while the other options stay admissible. \hfill$\square$
